\section*{ВВЕДЕНИЕ}
\addcontentsline{toc}{section}{ВВЕДЕНИЕ}
Виртуальная файловая система (ВФС) представляет собой программный слой абстракции, который предоставляет унифицированный интерфейс для доступа к разнородным физическим и логическим хранилищам данных.

В отличие от традиционных файловых систем, жестко привязанных к конкретному носителю или формату, ВФС оперирует не реальными секторами диска, а логическими объектами, отображая их для пользователя в виде привычной иерархии каталогов и файлов.

Ключевая особенность такой системы заключается в том, что источники данных могут быть географически распределены, иметь различную внутреннюю структуру и управляться независимо друг от друга.

Разработка собственной ВФС позволяет эмулировать поведение стандартных файловых операций -- создание, чтение, запись, перемещение и удаление -- над объектами, которые физически могут храниться, например, в записях базы данных или в оперативной памяти.

\emph{Цель настоящей работы} -- создание виртуальной файловой системы, а также изучение структуры FAT32 и CDFS. Для достижения поставленной цели необходимо решить \emph{следующие задачи:}
\begin{itemize}
\item изучение структуры FAT32 и CDFS;
\item проектирование виртуальной файловой системы;
\item проектирование интерфейса виртуальной файловой системы;
\item реализация виртуальной файловой системы и её интерфейса;
\item реализация файловых систем FAT32 и CDFS;
\item реализация командного интерпретатора.
\end{itemize}

\emph{Структура и объем работы.} Отчет состоит из введения, 4 разделов основной части, заключения, списка использованных источников, 2 приложений. Текст выпускной квалификационной работы равен \formbytotal{lastpage}{страниц}{е}{ам}{ам}.

\emph{Во введении} сформулирована цель работы, поставлены задачи разработки, описана структура работы, приведено краткое содержание каждого из разделов.

\emph{В первом разделе} на стадии описания технической характеристики предметной области приводится сбор информации о файловых системах.

\emph{Во втором разделе} на стадии технического задания приводятся требования к разрабатываемой виртуальной файловой системе.

\emph{В третьем разделе} на стадии технического проектирования представлены проектные решения для виртуальной файловой системы.

\emph{В четвертом разделе} приводится список функций, классов и их методов, использованных при разработке виртуальной файловой системы, производится тестирование разработанной виртуальной файловой системы.

В заключении излагаются основные результаты работы, полученные в ходе разработки.

В приложении А представлен графический материал.
В приложении Б представлены фрагменты исходного кода. 
